\documentclass[11pt,a4paper]{article}
\usepackage[T1]{fontenc}
\usepackage[margin=25mm]{geometry}
\usepackage{lmodern,microtype,amsmath,amssymb,booktabs,array}
\usepackage[hidelinks]{hyperref}
\setlength{\emergencystretch}{2em}
\newcommand{\SE}{\mathrm{SE}(3)}
\newcommand{\SO}{\mathrm{SO}(3)}
\newcommand{\diag}{\operatorname{diag}}
\newcommand{\norm}[1]{\left\lVert #1\right\rVert}
\title{Assessing Recovery Reliability in LiDAR-Inertial Odometry\\
\large Verified methods and protocol draft}
\author{Harshil Patel \and Daksha Lingampeta \and Nisarg Vyas}
\date{9 October 2026}
\begin{document}
\maketitle
\begin{abstract}
This document specifies an exit-conditioned evaluation of geometric health
signals from two existing LiDAR-inertial estimators. The evaluation separates
measurement geometry, short-window motion error, and accumulated drift. A
causal health decision is compared with a sustained local-motion accuracy label
computed offline from reference poses. The protocol fixes timing, availability,
reset handling, development selection, and geometry-level uncertainty before
unseen-layout exposure. Controlled simulation provides the quantitative
reference; a qualified handheld recording remains supporting evidence.
This is an in-progress methods document. Point-LIO development replication,
implementation freeze, held-out evaluation, and final manuscript conclusions
remain unfinished. No final performance or deployment claim is made here.
\end{abstract}

\section{Question and scope}
After a robot leaves weakly constrained geometry, when can an online geometric
health signal justify a decision that its local motion estimate is accurate?
Informative geometry and an accurate estimate are different measured objects.
The study evaluates existing estimators and diagnostic adaptations under one
fixed outcome protocol. Estimator inputs are LiDAR and IMU; reference poses
are confined to offline evaluation.

FAST-LIO2 registers raw points to an incremental map~\cite{fastlio}.
Point-LIO processes point-time measurements~\cite{pointlio}. We use pinned
implementations of these systems with read-only diagnostic exports. The
study's Point-LIO output is configured at frame end for the common decision
grid; this does not test its maximum published output bandwidth.

Prior work already includes a tunnel-exit recovery detector in LOFF~\cite{loff},
directional degeneracy characterization in DCReg~\cite{dcreg}, and estimator
health assessment against failure/error evidence~\cite{spectral}. The candidate
distinction is the prespecified exit-conditioned accuracy and timing
comparison. A final claim review must assess the actual results against this
overlap. No claim of being the first health or recovery detector is justified.

\section{Three quantities kept separate}
Let $T^r_{GB}(t)$ map comparison-body coordinates into the reference world and
$T^e_{WB}(t)$ map the same body into the estimator world. Quaternions are stored
as $(q_x,q_y,q_z,q_w)$; distance, time, and stored angle units are metres,
seconds, and radians.

\subsection{Local motion error}
For associated endpoints $t_0,t_1$, define
\begin{align}
 \Delta_r &= T^r_{GB}(t_0)^{-1}T^r_{GB}(t_1), &
 \Delta_e &= T^e_{WB}(t_0)^{-1}T^e_{WB}(t_1),\\
 E_{\rm rel} &= \Delta_r^{-1}\Delta_e
              =\begin{bmatrix}R_{\rm rel}&p_{\rm rel}\\0&1\end{bmatrix}.
\end{align}
The translation error is $\norm{p_{\rm rel}}_2$ and rotation error is
$\arccos\!\left(\operatorname{clip}((\operatorname{tr}R_{\rm rel}-1)/2,-1,1)\right)$.
Rates divide by the actual endpoint separation. Local error requires no
global alignment and no post-event re-alignment.

Primary labels use non-overlapping one-second windows on the integer-second
simulation grid. Each nominal boundary is associated once to the nearest valid
pose within $\pm0.05$~s, with ties assigned to the earlier pose. Adjacent windows
reuse the same shared boundary pose. Valid durations lie in $[0.9,1.1]$~s;
windows crossing a reset, missing a boundary, or lacking reference coverage are
unavailable. Three-second translation/rotation rates are secondary sensitivity
outcomes. Unavailable errors never become zero-valued observations.

\subsection{Accumulated drift}
At the declared pre-entry anchor $t_a$, use one rigid, unit-scale alignment
\begin{equation}
 T_{GW}=T^r_{GB}(t_a)T^e_{WB}(t_a)^{-1},\qquad
 E_{\rm acc}(t)=T^r_{GB}(t)^{-1}T_{GW}T^e_{WB}(t).
\end{equation}
Accumulated error uses $E_{\rm acc}(t)$ throughout the run. The alignment is not
refitted after entry, exit, or reset. This quantity is reported separately
from local recovery and does not define a false-healthy event.

\section{Measurement-geometry diagnostics}
Both backends expose their own accepted point-to-plane measurement rows.
The geometry diagnostic excludes estimator priors and is not a posterior
covariance. With the six columns ordered as translation then rotation,
let $H=\sum_i j_i^\top j_i$ and
\begin{equation}
 S_L=\diag(1,1,1,L^{-1},L^{-1},L^{-1}),\qquad
 G_L=\frac{S_L H S_L}{N\,0.001}.
\end{equation}
The main scalar is $\lambda_{\min}(G_{3\,\mathrm m})$; $L=1$ and $5$~m are
retained sensitivity outputs. The normalization is fixed by the protocol.
In particular, it does not assert that Point-LIO applies a measurement noise
weight of $0.001$ in its filter.

FAST-LIO uses the last accepted scan linearization. Point-LIO requires an
explicit adaptation because its state changes between point-time groups.
For group $g$, retain the rotation $R_g$ at the same pre-update linearization
that produced the accepted Jacobian. Transport rotational columns to a common
world tangent basis before pooling:
\begin{equation}
 j_{{\rm common},g}=[j_{t,g},\;j_{r,g}R_g^\top],\qquad
 H_{\rm frame}=\sum_g j_{{\rm common},g}^\top j_{{\rm common},g}.
\end{equation}
This is a frame-pooled geometry score rather than information for a single
scan-end posterior. Point-LIO correspondences, poses, and fitted operating
point remain distinct from FAST-LIO's.

For the DCReg detector adaptation, permute the raw Hessian into rotation then
translation order. Its Schur blocks are
\begin{align}
 S_R &= H_{RR}-H_{Rt}H_{tt}^{\dagger}H_{tR},\\
 S_t &= H_{tt}-H_{tR}H_{RR}^{\dagger}H_{Rt}.
\end{align}
The pseudoinverse uses relative rank tolerance $10^{-12}$. Subspace matrices
are symmetrized and numerically checked before spectral analysis. Direction
ratios are $\kappa_i=\lambda_{\max}/\lambda_i$; a ratio strictly greater than
10 is flagged. Zero/rank-deficient directions have explicit unbounded ratios.
This evaluates the directional detector on backend-specific rows, without
running DCReg's complete registration or mitigation system.

At least six accepted rows and a finite numerically valid aggregate are
required. Ordinary empty groups do not invalidate a completed frame whose
total accepted count is sufficient. Missing groups, absent frame records,
crashes, and unfinished frames are incomplete execution. Startup, empty input,
insufficient correspondences, numeric failure, and resets remain explicit
unavailable states. A finite rank-deficient matrix remains valid geometry
evidence, including a zero smallest eigenvalue.

\section{Recovery labels and causal decisions}
Geometry exit is annotated independently of estimator error and indicators.
In simulation it is the analytical route crossing of the exit plane. Recovery
is the first post-exit triple of consecutive non-overlapping one-second
windows whose translation errors are at most $0.20$~m and rotation errors at
most $5^\circ$. All triple endpoints must lie by $t_{\rm exit}+20$~s.
Onset is the first window's actual start; confirmation is the third window's
actual end. Later failure of a required window is a relapse. Completed events
without a qualifying triple are non-recoveries; crashed/incomplete attempts
are not relabelled as non-recovery.

The recovery term is this operational post-exit accuracy label. It does not
by itself establish that local accuracy had previously been lost in every
corridor event. That distinction must remain visible in interpretation.

At each 10~Hz decision tick, use only the latest diagnostic at or before that
tick, no more than $0.2$~s old and in the same estimator segment. An invalid
latest record cannot be skipped to reuse an older healthy record. Three
distinct consecutive healthy source records confirm an alarm. Missing,
unhealthy, stale, or cross-segment evidence clears confirmation; repeated
use of one source record does not count as an additional sample.

\begin{table}[ht]
\centering\small
\begin{tabular}{p{0.24\linewidth}p{0.68\linewidth}}
\toprule Outcome & Definition and denominator\\\midrule
False-healthy event & Confirmed health active during the initial unrecovered
interval; divide by truth-eligible corridor events with an available decision
in that interval. Also report the unconditional rate and availability.\\
Recovery sensitivity & Post-onset confirmation before recovery ends, divided
by all truth-recovered eligible events. Missing scores or absent qualifying
transitions are misses.\\
Detection delay & Confirmation time minus true onset for qualifying alarms;
report detection fraction and median among detections, with other cases censored.\\
Post-relapse health & Report continued confirmed health after relapse with
conditional and unconditional denominators.\\
\bottomrule
\end{tabular}
\caption{Frozen outcome meanings. Empty denominators are undefined, not zero.}
\end{table}

\section{Development, unseen scenes, and uncertainty}
Development uses 32 geometry seeds (14--45), each with a corridor and matched
nondegenerate control. Paired IMU streams are identical. Four prescribed
first/last same-input backend repeats add no independent geometry events.
Select each backend's scalar operating point separately using only its formal
corridor development events: exhaust unique finite valid scores plus an
explicit no-alarm candidate; maximize recovery sensitivity subject to at most
10\% conditional false-healthy events, breaking ties toward the higher
threshold. No feasible positive-sensitivity choice is reported explicitly.
Controls and the isolated feasibility replay do not train an operating point.

The current final plan contains 48 paired unseen layouts, 12 in each of four
fixed length/width strata. Geometry-only eligibility screening proceeds in
reserved-seed order after implementation and replication freezes pass. Final
outcomes cannot tune methods, choose replacement scenes, or alter thresholds.
Transfer claims are limited to these unseen straight-corridor layouts.

Event proportions use 95\% Wilson intervals. Continuous paired outcomes first
reduce valid planned windows to one per-seed median, then average paired
differences within strata and weight the four stratum means equally. Use
10,000 stratified geometry-seed bootstrap replicates with PCG64 seed 20260926,
preserving corridor/control pairs and fixed stratum order. Percentiles use
NumPy's linear method. Undefined ROC/AP replicates make the corresponding
curve interval unavailable. Incomplete final pairs after the one technical
retry remain failures; successful-run summaries then remain descriptive
rather than confirmatory evidence at the planned sample size.

\section{Reproduction and evidence limits}
Record input, configuration, source, native executable, environment, output,
and failure identities. Cache only verified matching upstream results. Keep
interrupted attempts and technical retries separate from independent scenes.
Reference truth is never given to the estimator or diagnostic producer.
Implementation checks include analytic matrices, reset barriers, frame/group
coverage, numeric unavailability, pose identity joins, output-flush checks,
and unchanged estimator output with diagnostics enabled.

Simulation supplies ideal analytical reference with explicit sensor timing,
visibility, noise, and IMU assumptions. It does not establish physical drone
performance. The Hilti-Oxford handheld example has sparse usable post-exit
reference and partly map-registration-derived truth; it cannot support the
primary sustained-recovery conclusion. The publication distinction and final
findings require independent review after untouched evaluation.

\begin{thebibliography}{9}
\bibitem{fastlio} W. Xu, Y. Cai, D. He, J. Lin, and F. Zhang,
``FAST-LIO2: Fast Direct LiDAR-inertial Odometry,''
\href{https://arxiv.org/abs/2107.06829}{arXiv:2107.06829}, 2021.
\bibitem{pointlio} D. He, W. Xu, N. Chen, F. Kong, C. Yuan, and F. Zhang,
``Point-LIO: Robust High-Bandwidth Light Detection and Ranging Inertial
Odometry,'' \emph{Advanced Intelligent Systems}, 2023,
\href{https://doi.org/10.1002/aisy.202200459}{doi:10.1002/aisy.202200459}.
\bibitem{loff} ``LOFF: LiDAR and Optical Flow Fusion Odometry,''
\emph{Drones}, vol. 8, no. 8, article 411, 2024,
\href{https://doi.org/10.3390/drones8080411}{doi:10.3390/drones8080411}.
\bibitem{dcreg} ``DCReg: Decoupled Characterization for Efficient Degenerate
LiDAR Registration,'' \href{https://arxiv.org/abs/2509.06285}{arXiv:2509.06285},
version 3, 2026.
\bibitem{spectral} C. Lanegger, H. Oleynikova, R. Siegwart, and M. Pantic,
``When Your State Estimator Has Lost The Plot: Detecting Estimator Failures
Via Spectral Analysis,'' \href{https://arxiv.org/abs/2608.10623}{arXiv:2608.10623}, 2026.
\end{thebibliography}
\end{document}
